\section{Why this game is worth studying}
\label{sec:domain}

Fish sits at an unusual point in the space of imperfect-information games, and the point is
unusual in a way that makes certain questions cheap to ask that are expensive everywhere else. This
section says what those questions are, because they are the reason the rest of the paper exists.

\subsection{The hidden state is the deal, and the posterior over it is exact}
\label{sec:domain-exact}

Every card movement in Fish is public. A card changes hands only through a successful ask, and the
ask, its target and its outcome are seen by everyone; a declaration reveals an entire half-suit's
allocation whether it succeeds or fails. It follows that the \emph{only} thing any player does not
know is the initial deal, and that the public record determines every subsequent card position given
that deal.

This is a stronger property than it sounds. In most hidden-hand card games the reachable-history set
is what must be filtered, and history filtering is FNP-complete in general
\cite{solinas-history}. Fish escapes: there is exactly one object to reconstruct, and the
constraints on it are logical rather than probabilistic. A player who has been told ``I do not hold
the seven of hearts'' knows it, forever, with certainty. The observation history is a conjunction of
hard constraints on a fixed unknown, and the posterior over deals consistent with those constraints
is computable in closed form (\S\ref{sec:engine-belief}).

The consequence for research is that Fish lets you \emph{separate} inference from policy. In poker
or Bridge, a claim that a policy is weak can always be answered with ``its beliefs are wrong''. Here
the beliefs can be made provably right, at which point every remaining loss is attributable to the
policy. That separation is what makes the negative results of \S\ref{sec:arc} and
\S\ref{sec:rejected} meaningful rather than merely discouraging.

It also sets up the single most surprising finding in this corpus, which
\S\ref{sec:arc-vsix-belief} reports: making the belief exact makes the policy \emph{worse}.

\subsection{The team is the unit, and the team cannot talk}
\label{sec:domain-team}

Three teammates share a payoff, hold different private information, and have no communication
channel beyond the public actions themselves. That is precisely the decentralised control problem
the common-information approach formalises: reformulate from the point of view of a coordinator who
knows the common information and issues \emph{prescriptions}, maps from each controller's private
information to its action \cite{nayyar2013common}. In Fish the common information is the entire
public history, so the coordinator's information state is exactly the belief object the engine
already maintains.

Two things follow that matter for evaluation rather than for algorithms. First, the coordinator's
problem is not solvable here --- a prescription maps every possible nine-card hand to an ask --- so
the reduction is a correctness lens and not a method. Second, and more usefully, \emph{any
deterministic procedure that reads only public state and its own hand is automatically common
knowledge among the three teammates}, with no shared seed and no communication. That is a cheap
version of a property the cooperative-search literature has to work for \cite{lerer2020sparta}, and
it is why the side-channel definition of \S\ref{sec:threat} can be made mechanical.

The qualifier is load-bearing, and \S\ref{sec:arc-vsix-ties} is where it bites: a randomised
tie-break loses the property unless its randomisation is seeded from the public history.

\subsection{A team of published, deterministic, identical agents inverts a standard assumption}
\label{sec:domain-inverted}

The team-correlation literature uniformly assumes the correlation device is the team's own ---
private, and unknown to the opponent \cite{celli-gatti,farina-collusion,zhang-tbdag}. A homogeneous
team of three \emph{identical published deterministic} agents inverts that: the device is common
knowledge, and an adversary may read it. The threat model of \S\ref{sec:threat} takes that
seriously, and \S\ref{sec:exploit} reports what an adversary that reads the target's own policy and
inverts its transcript actually achieves.

We found no prior work covering this case, and \S\ref{sec:related} states the scope of that search
rather than claiming priority beyond it.

\subsection{The domain is cheap, and that is a methodological asset}
\label{sec:domain-cheap}

A Fish deal plays out in roughly a hundred public events. The engine used here plays hundreds of
complete six-player games per second per core with exact inference running in every seat. This is
not a boast about engineering; it is what makes the evaluation programme in \S\ref{sec:methodology}
possible at all. A study that needs 48{,}000 paired games to resolve a three-point effect, two
disjoint holdout banks for every claim, a 31-member opponent panel scored identically for four arms,
and eight independent adversary searches, can run all of that on one desktop in a day. In a domain
where a single evaluation cell costs GPU-hours, the honest version of this paper could not have been
written.

The corresponding liability is stated in \S\ref{sec:limitations}: cheap evaluation is not external
evaluation, and every number here was produced by one engine that this project wrote.

\subsection{What the game does \emph{not} offer}
\label{sec:domain-not}

Two disclaimers, because both bear on how the results should be read.

Fish's perfect-information game is degenerate, not merely fused. A clairvoyant player names only
cards the target actually holds, so never fails an ask, never loses the turn, and declares each
half-suit as soon as their team holds all of it. Almost every reachable position carries the same
perfect-information value. Determinization --- the standard recipe for hidden-hand card games
\cite{levy,ginsberg-gib} --- therefore collapses here to choosing the ask most likely to hit, blind
to the information an ask leaks, the information a refusal supplies, and the option value of
postponing a declaration. \S\ref{sec:related-determinization} places this against the literature,
and \S\ref{sec:arc-vsix-search} reports what the project measured when it built such a search
anyway.

And Fish has no passive continuation. A player must ask or declare; there is no equivalent of poker's
check/call rollout, which is what the standard cheap exploitability probe is built on
\cite{lisy2017lbr}. \S\ref{sec:towardone} names this as the single largest obstacle between this
corpus and a defensible near-optimality claim.
